\documentclass[10pt,twocolumn,letterpaper]{article}

\usepackage{cvpr}
\usepackage{times}
\usepackage{epsfig}
\usepackage{graphicx}
\usepackage{amsmath}
\usepackage{amssymb}
\usepackage{booktabs}
\usepackage{multirow}
\usepackage{array}

\usepackage[pagebackref=true,breaklinks=true,letterpaper=true,colorlinks,bookmarks=false]{hyperref}

\def\cvprPaperID{****}
\def\httilde{\mbox{\tt\raisebox{-.5ex}{\symbol{126}}}}

\ifcvprfinal\pagestyle{empty}\fi
\begin{document}

%%%%%%%%% TITLE
\title{What Happens Next? Action Anticipation on a 33-class subset of Something-Something V2}

\author{Khalil Ayadi\\
\'Ecole Polytechnique --- CSC\_43M04\_EP\\
{\tt\small khalil.ayadi@polytechnique.edu}
}

\maketitle

%%%%%%%%% ABSTRACT
\begin{abstract}
We address the \emph{What Happens Next?} challenge, a 33-class action
anticipation task built from the early frames of Something-Something V2
(SSv2) clips. Two tracks are considered: a closed-world track (training
from scratch, no external data or weights) and an open-world track
(pretrained models allowed). For the closed track, we benchmark a
CNN--LSTM with attention pooling, a two-stream spatio-temporal
transformer, and the Temporal Shift Module (TSM). We observe that
\emph{architectures that encode temporal reasoning through cheap,
structural priors} (TSM channel shift, two-stream motion differencing)
clearly outperform pure-attention models trained from scratch on this
small dataset. For the open track, we adapt VideoMAE, TimeSformer,
MViTv2 and V-JEPA~2 (ViT-L 0.4B and ViT-g 1B) through HuggingFace
wrappers, and add a LoRA fine-tuning path to keep the 1B variant
tractable. Our strongest submission is a late ensemble of heterogeneous
checkpoints whose weights and per-model temperatures are learned by
\emph{simplex projected SGD}, calibrated on leakage-free K-fold
out-of-fold predictions, and refined at inference with 5-crop
test-time augmentation. We also document two inference-time corrections
(Bayesian log-prior subtraction, temperature calibration). Source code
and detailed experimental logs are provided alongside this report.
\end{abstract}

%-------------------------------------------------------------------------
\section{Problem description}

The challenge is a 33-class subset of Something-Something V2~\cite{Goyal17}.
Each clip is represented as four extracted frames covering the
\emph{beginning} of an action; the system must predict the full action
class (e.g.~\textit{Pretending to throw something}, \textit{Pouring
something into something}). The metric is top-1 accuracy on a held-out
test set, with top-5 reported as a secondary metric.

Two factors make this dataset specifically hard:
\begin{itemize}
    \item \textbf{Partial observation.} Only 4 frames are available, and
    several classes are temporal twins of more common classes
    (e.g.~\textit{Picking up} vs. \textit{Pretending to pick up}) that
    become visually identical in the first few frames.
    \item \textbf{Class imbalance.} Train counts range from
    $\sim$$160$ for \textit{Spilling} to $\sim$$2{,}250$ for
    \textit{Throwing}, and class index $27$ is entirely missing from the
    training split.
\end{itemize}

The challenge defines two tracks: closed world (Track A) forbids any
pretrained weights or external data, and open world (Track B) allows
them. We participated in both. Our code is built on top of the
instructor-provided Hydra+PyTorch skeleton; we only added new model
files, new YAML configs, and \texttt{build\_model} branches, plus
opt-in flags in the training loop. Inference scripts
(\texttt{evaluate.py}, \texttt{create\_submission.py}) were never
modified.

%-------------------------------------------------------------------------
\section{Method}

\subsection{Pipeline and shared training recipe}

Every model accepts a $(B, T, C, H, W)$ tensor of frames and outputs
$(B, K)$ logits with $K=33$. With $T=4$ pinned by the dataset, the
design problem reduces to: \emph{how do we extract spatial features per
frame, and how do we mix them across the 4 time steps?} The dataset
loader, optimizer (AdamW, which beats the canonical SGD recipe on this
task, Table~\ref{tab:opt}), AMP, augmentation knobs, label smoothing,
and class-balanced sampling are shared by all closed-track models.
Each architecture only chooses the spatial encoder and the temporal
aggregation.

\subsection{Data augmentation}
\label{sec:aug}

All augmentation knobs are exposed through Hydra so individual
ablations are one-line CLI overrides. For from-scratch training
(Track~A) we ablated each transform on the validation set; the outcome
guided a deliberately \emph{minimal} final recipe, contrary to our
initial intuition that more augmentation would help on a small dataset.

\begin{itemize}
    \item \textbf{Kept.} Random resized crop (scale $[0.6, 1.0]$) and
    ColorJitter (strength $0.2$): both consistently improved or left
    top-1 unchanged, and add no label noise.
    \item \textbf{Removed after ablation.} Horizontal flip, small
    rotation, sharpness and Gaussian blur each \emph{degraded} top-1.
    Our best from-scratch checkpoints are exactly those with all four
    disabled (Appendix~A).
    \item \textbf{Mixing.} Mixup~\cite{Zhang18mixup} ($\alpha=0.2$) and
    CutMix~\cite{Yun19cutmix} ($\alpha=1.0$), evaluated both on and off
    as a regularizer; kept as a diversity lever across ensemble members.
\end{itemize}

A subtle but important point explains the flip result:
\emph{horizontal flip is label-breaking} on SSv2, which contains
direction-sensitive classes (\textit{Pushing something from left to
right} vs.~\textit{\dots from right to left}). For the two-stream
model, flipping additionally corrupts the motion stream, which encodes
the \emph{sign} of horizontal motion. This is also why the test-time
augmentation of Sec.~\ref{sec:tta} excludes the flip.

\subsection{Closed-track models (Track~A)}
\label{sec:closed}

We benchmarked three architectures.

\paragraph{CNN--LSTM.} ResNet-18 features per frame, LayerNorm and
feature dropout, a bidirectional 2-layer LSTM with hidden $256$,
attention pooling over the $T=4$ outputs, and a dropout-$0.5$ head.
Orthogonal init on $W_{hh}$, Xavier on $W_{ih}$, forget-gate bias
initialised to $1.0$~\cite{Jozefowicz15}. We deliberately replaced
``take the last LSTM hidden'' with additive attention, since on $T=4$
the last hidden state discards $75\%$ of the temporal information.

\paragraph{Two-stream spatio-temporal transformer.} Two parallel
shallow ResNet backbones (RGB and frame-difference), each topped by a
single spatial self-attention layer over $H\times W$ positions. A small
temporal transformer aggregates the per-stream, per-frame tokens through
a learnable \texttt{[CLS]} token. We empirically reverted from
mean-pooling back to \texttt{[CLS]} aggregation after observing val
accuracy collapse with mean pooling (see Sec.~\ref{sec:failures}).

\paragraph{Temporal Shift Module (TSM).} A 2D ResNet variant where, at
every residual block, $1/8$ of channels are shifted forward and $1/8$
backward along the temporal axis~\cite{Lin19tsm}. The shift carries
zero parameters and negligible FLOPs but gives the 2D CNN access to
neighbouring frames inside every block. We expose the backbone depth as
a flag (ResNet-18 by default, ResNet-34 layout $(3,4,6,3)$ optional),
and the temporal aggregation as \texttt{mean}, attention or a depthwise
1-D \texttt{tconv}.

\subsection{Open-track models (Track~B)}

For Track~B, we added HuggingFace wrappers (under the \texttt{b\_}
prefix) for four SSv2-pretrained backbones:

\begin{itemize}
    \item \textbf{VideoMAE base} ($\sim$$85$M params,
    \cite{Tong22videomae}): masked-autoencoding pre-training on SSv2,
    then supervised fine-tuning.
    \item \textbf{TimeSformer base}~\cite{Bertasius21timesformer}:
    \emph{divided space-time attention} ViT, fine-tuned on SSv2.
    \item \textbf{MViTv2}~\cite{Li22mvitv2}: multiscale ViT with
    stage-wise spatial pooling, an architecturally diverse complement
    to the plain ViT family above.
    \item \textbf{V-JEPA~2} (Meta, 2025): self-supervised video
    representation learning. We integrated both the ViT-L
    ($\sim$$400$M) and ViT-g ($\sim$$1.1$B) variants.
\end{itemize}

For the V-JEPA~2 1B model, full fine-tuning is fragile and expensive,
so we added a \textbf{LoRA path} using HuggingFace
\texttt{peft}~\cite{Hu22lora}. Rank-$16$ LoRA adapters are inserted into
the attention projections (\texttt{q/k/v/o\_proj}), and only the new
classifier head is kept fully trainable. This freezes
$\sim$$99.5\%$ of the parameters, makes training 2--3$\times$ faster,
and--crucially--makes catastrophic forgetting of the SSv2 prior
impossible.

\subsection{Inference-time corrections}

Three corrections are applied at inference, model-agnostically.

\paragraph{Log-prior subtraction.} The training set imposes an implicit
prior $\log P(y_{\text{train}})$ on the network output. For pairs of
visually-similar classes, the model defaults to whichever class has more
training samples. We subtract $\alpha \cdot \log P(y_{\text{train}})$
from each logit at inference~\cite{Menon20logitadj}, with the missing
class~$27$ crushed to $-\infty$. We found $\alpha=0.5$ a safe sweet spot;
on a test set that matches the train distribution, full subtraction
($\alpha=1$) over-corrects and lowers top-1, so we keep it optional.

\paragraph{Per-model temperature scaling.} Before averaging softmaxes
in the ensemble, we fit a per-model temperature $T_i$ on the NLL of the
calibration set~\cite{Guo17calibration}. This does not change a model's
top-1 but rescales over-confident models so that they do not dominate
the average.

\paragraph{Test-time augmentation (TTA).}
\label{sec:tta}
At inference we average the softmax over five spatial crops per clip
(center $+$ four corners, after resizing the short side at the classic
$256/224$ ratio). Every frame of a clip receives the \emph{same} crop so
temporal motion is preserved, and we deliberately exclude the
horizontal flip of standard 10-crop TTA (Sec.~\ref{sec:aug}). TTA is
applied identically to the calibration and the test set, and yields a
consistent $+0.5$--$1.5$ top-1 points for no retraining.

\subsection{Ensemble: simplex-SGD weights with leakage-free calibration}
\label{sec:ensemble}

Our final submission is a weighted average of per-model
temperature-scaled softmaxes,
\[
P_{\text{ens}}(y\mid x) = \sum_i w_i\,
\text{softmax}\!\left(\frac{\ell_i(x)}{T_i}\right),
\]
where $\ell_i$ are the raw logits of model $i$. Rather than a heuristic
accuracy-power weighting, we \emph{learn} the weights and temperatures
directly. We parameterize $w = \text{softmax}(\theta)$, which keeps the
weights on the probability simplex by construction ($w_i\!\ge\!0$,
$\sum_i w_i\!=\!1$), and optionally a per-model log-temperature, then
minimize the ensemble NLL on a calibration set by Adam (we call this
\emph{simplex-SGD}). A small $\ell_2$ penalty on $\theta$ keeps the
solution close to uniform on noisy calibration data. Because the weights
cannot go negative, adding a weak-but-diverse model is low-risk: the
optimizer simply drives its weight toward zero rather than letting it
harm the average.

\paragraph{Leakage-free calibration (K-fold OOF).} Fitting the weights
on data a model was trained on would reward over-confident memorization.
We therefore calibrate either on the held-out \texttt{val/} directory
(never seen at training), or on \emph{K-fold out-of-fold} (OOF)
predictions: each family is trained $K=5$ times, each fold holding out a
disjoint $20\%$ of the train set, so that stitching the five held-out
predictions yields one honest softmax per training clip. The simplex
weights fit on these OOF predictions transfer to the test set without
leakage. The script verifies that the union of the five held-out folds
covers the train set exactly (no gap, no duplicate) before stacking.

%-------------------------------------------------------------------------
\section{Implementation}

Every model lives in \texttt{src/models/*.py}, with one YAML in
\texttt{configs/model/} and one ``experiment'' YAML in
\texttt{configs/experiment/} overriding the training defaults. A K-fold
TSM family is launched with one CLI per fold:

\begin{verbatim}
python src/train.py experiment=tsm \
  +dataset.kfold_num=5 +dataset.kfold_index=0 \
  training.epochs=30 model.temporal_pool=mean \
  training.use_class_balanced_sampler=false \
  training.use_color_jitter=true \
  training.use_horizontal_flip=false \
  training.checkpoint_path=tsm_fold0.pt
\end{verbatim}

\texttt{oof\_predict.py} then stitches the five fold checkpoints into one
leakage-free softmax tensor, and \texttt{ensemble\_predict.py} takes an
arbitrary list of checkpoints plus a calibration source (OOF files or a
labelled directory), fits the simplex-SGD weights, optionally applies
TTA, and writes the submission CSV. Inference scripts reload each
architecture from the \emph{config merged into the checkpoint}, so a
checkpoint produced by any branch reloads through a single
\texttt{build\_model} call.

\paragraph{What is ours vs.~provided.} The instructor-provided skeleton
consists of the dataset loader, the training loop scaffold, the
inference scripts, and an initial \texttt{cnn\_lstm} baseline.
\emph{Our additions} are: the two-stream transformer and TSM models
(with the ResNet-34 depth flag); the augmentation toggles in
\texttt{utils.build\_transforms}; mixup/cutmix/class-balanced sampling;
the focal loss and logit-adjustment training paths; the log-prior
calibration; the K-fold split and \texttt{oof\_predict.py}; the
simplex-SGD$+$TTA ensemble in \texttt{ensemble\_predict.py}; and the
LoRA-aware V-JEPA~2 wrapper. The HuggingFace pretrained checkpoints used
for Track~B are external (cited above); we re-use the standard
\texttt{peft} library for LoRA.

\paragraph{Engineering encountered.} Three issues consumed significant
time and are documented in \texttt{problems.md}: (1) a DataLoader
bottleneck (GPU oscillating $0\%\!\rightarrow\!100\%$) fixed with AMP,
TF32, \texttt{cudnn.benchmark}, non-blocking transfers and a higher
\texttt{prefetch\_factor}; (2) an OOM-kill of the workers fixed by
lowering \texttt{prefetch\_factor} and defaulting
\texttt{persistent\_workers} to false; (3) a prior-calibration bug where
the missing class $27$ received a $+20.7$~nat boost and collapsed top-1
to zero, fixed by assigning $-\infty$ to absent classes.

%-------------------------------------------------------------------------
\section{Results}

\subsection{Closed track (Track~A)}

Table~\ref{tab:closed} reports validation top-1 / top-5 of the main
closed-track runs (held-out \texttt{val/} directory, raw logits, best
checkpoint per run). The full recipe ablation is in Appendix~A.

\begin{table}[h]
\centering
\small
\begin{tabular}{lcc}
\toprule
\textbf{Model} & \textbf{Top-1} & \textbf{Top-5} \\
\midrule
CNN--LSTM (BiLSTM + attn pool)   & 0.18 & 0.52 \\
Two-stream (CLS, $bc{=}32$)      & 0.30 & 0.63 \\
TSM (ResNet-18, shift $1/8$)     & 0.33 & 0.63 \\
\midrule
\textit{Ensemble (TSM + two-stream)} & \textbf{0.40} & \textbf{0.73} \\
\bottomrule
\end{tabular}
\caption{Closed-track validation top-1 / top-5. The ensemble uses
simplex-SGD weights with per-model temperature calibration and 5-crop
TTA (Sec.~\ref{sec:ensemble}--\ref{sec:tta}).}
\label{tab:closed}
\end{table}

\subsection{Open track (Track~B)}

Table~\ref{tab:open} reports the same metrics on the Track~B models.
The strongest single model is V-JEPA~2 fine-tuned via LoRA.

\begin{table}[h]
\centering
\small
\begin{tabular}{lcc}
\toprule
\textbf{Model} & \textbf{Top-1} & \textbf{Top-5} \\
\midrule
TimeSformer base (full FT)            & 0.41 & 0.74 \\
VideoMAE base (full FT)               & 0.46 & 0.78 \\
MViTv2 base (full FT)                 & 0.45 & 0.77 \\
V-JEPA 2 ViT-L 0.4B (LoRA r=16)       & 0.53 & 0.82 \\
V-JEPA 2 ViT-g 1B (LoRA r=16)         & \textbf{0.58} & \textbf{0.85} \\
\bottomrule
\end{tabular}
\caption{Open-track validation top-1 / top-5 (indicative of typical
runs; minor variation across seeds).}
\label{tab:open}
\end{table}

\subsection{What worked, what did not}
\label{sec:failures}

\paragraph{Why TSM and the two-stream model work on $T=4$.}
Both bake a strong \emph{temporal inductive bias} into the architecture
without adding parameters: TSM shifts a fraction of channels along time
at every residual block, and the two-stream model precomputes
$\Delta_t = \text{frame}_t - \text{frame}_{t-1}$ as an explicit motion
channel. With only four frames, every architectural ``hint'' that the
temporal axis exists is worth more than fitted attention weights.

\paragraph{Capacity from scratch.}
On this small dataset, wider/deeper is not better: the two-stream model
with $bc{=}64$ loses $6$--$8$ top-1 points to $bc{=}32$
(Appendix~A), and a ResNet-34 TSM tends to underperform ResNet-18 in
isolation. We therefore keep the narrow variants as the leaders and use
the larger ones only as decorrelated ensemble members.

\paragraph{Where calibration helps.}
Per-class recall on \texttt{val} is worst on the \emph{semantic-twin}
classes (\emph{Picking up}, \emph{Pretending to put into},
\emph{Pretending to throw}, \emph{Spilling}). Log-prior calibration with
$\alpha=0.5$ recovers $5$--$15$ recall points on these, at the cost of
$1$--$3$ top-1 points overall. For the ensemble, per-model temperature
was the single most impactful pre-processing step.

\paragraph{Failure case analysis.}
The dominant residual error remains the \emph{pretending-vs-real}
confusion when the first 4 frames show no contact between hand and
object. This is fundamentally \emph{partial observability}: the
disambiguating frame is by construction not in the input. The dataset
itself sets an irreducible Bayes error on these classes.

%-------------------------------------------------------------------------
\section{Conclusion and future work}

We benchmarked three closed-world architectures (CNN--LSTM, two-stream
transformer, TSM) and four open-world pretrained backbones on a
partial-observation, long-tail SSv2 subset, and combined them into a
calibrated ensemble. The two strongest \emph{conceptual} takeaways are:
(i) for tiny temporal contexts ($T=4$), architectural temporal bias
(TSM channel shift, two-stream motion) beats fitted attention from
scratch; and (ii) honest ensemble calibration --- simplex-SGD weights
fit on K-fold out-of-fold predictions, plus per-model temperature and
5-crop TTA --- is a cheap, high-leverage step that the heuristic
accuracy-power weighting does not match.

\paragraph{Extensions.} The most promising directions we did not have
time to pursue are: (1) increasing $T$ from $4$ to $8$ by re-extracting
frames closer to the action onset; (2) an \emph{adapter ensemble} where
each member is a separate LoRA on the same V-JEPA~2 backbone (cheap,
since adapters are tiny); (3) replacing the dense per-class log-prior
with a \emph{semantic-pair} prior that damps only the dominant twin of
each confused pair, leaving unrelated classes alone.

%-------------------------------------------------------------------------
\appendix
\section{Appendix --- Detailed experiments}

A full appendix of per-run logs accompanies the submission; below we
report only the ablations that drive the design choices above.

\subsection{TSM recipe ablation}

Table~\ref{tab:tsm-ablation} sweeps the class-balanced sampler
(\texttt{cbs}), the LR schedule (cosine vs.~flat) and the temporal
pooling (mean vs.~\texttt{tconv}) at fixed seed, with flip / rotation /
sharpness / blur disabled and crop $+$ color jitter on. Mean pooling,
cosine schedule and no class-balanced sampler dominate.

\begin{table}[h]
\centering
\small
\begin{tabular}{llcc}
\toprule
\textbf{cbs} & \textbf{sched / pool} & \textbf{Top-1} & \textbf{Top-5} \\
\midrule
off & cosine / mean   & 0.326 & 0.637 \\
off & cosine / tconv  & 0.320 & 0.619 \\
off & flat / mean     & 0.316 & 0.653 \\
off & flat / tconv    & 0.322 & 0.644 \\
on  & cosine / mean   & \textbf{0.329} & 0.632 \\
on  & cosine / tconv  & 0.309 & 0.621 \\
on  & flat / mean     & 0.308 & 0.633 \\
on  & flat / tconv    & 0.293 & 0.614 \\
\bottomrule
\end{tabular}
\caption{TSM ablation (val, ResNet-18, $30$ epochs). The mean/cosine
rows lead; the single best run uses the sampler, but the no-sampler
configuration wins $3/4$ pairings and is our default. Seed variance is
$\le 0.005$ top-1.}
\label{tab:tsm-ablation}
\end{table}

\subsection{Optimizer: AdamW vs.~SGD}

Table~\ref{tab:opt} pairs the AdamW and SGD recipes at otherwise
identical settings (ResNet-18, sampler on). AdamW (lr $10^{-3}$) wins
top-1 in every pairing against the canonical SGD recipe
(lr $10^{-2}$, momentum $0.9$) by $0.3$--$0.8$ points; top-5 is a wash.
AdamW is therefore our default optimizer for TSM.

\begin{table}[h]
\centering
\small
\begin{tabular}{lcc}
\toprule
\textbf{schedule / pool} & \textbf{SGD} & \textbf{AdamW} \\
\midrule
cosine / mean   & 0.329 & \textbf{0.332} \\
cosine / tconv  & 0.309 & \textbf{0.317} \\
flat / tconv    & 0.293 & \textbf{0.297} \\
\bottomrule
\end{tabular}
\caption{TSM optimizer comparison (val top-1). AdamW at $10^{-3}$ beats
the SGD recipe at $10^{-2}$ in all pairings; top-5 differences are
within noise.}
\label{tab:opt}
\end{table}

\subsection{Two-stream width and sampler}

Table~\ref{tab:twostream-ablation} shows the two-stream model is
strongly hurt by the wide $bc{=}64$ stem (overfit from scratch) and
mildly hurt by the class-balanced sampler; label smoothing is a wash.

\begin{table}[h]
\centering
\small
\begin{tabular}{lcccc}
\toprule
 & \multicolumn{2}{c}{$bc{=}32$} & \multicolumn{2}{c}{$bc{=}64$} \\
\textbf{config} & T-1 & T-5 & T-1 & T-5 \\
\midrule
$ls{=}0.1$, no cbs & \textbf{0.297} & 0.628 & 0.227 & 0.521 \\
$ls{=}0.1$, cbs    & 0.288 & 0.606 & 0.204 & 0.515 \\
$ls{=}0.0$, no cbs & 0.295 & 0.619 & 0.231 & 0.550 \\
$ls{=}0.0$, cbs    & 0.285 & 0.609 & 0.222 & 0.521 \\
\bottomrule
\end{tabular}
\caption{Two-stream ablation (val). $bc{=}32$ beats $bc{=}64$ by
$6$--$8$ points everywhere; \texttt{nocbs} beats \texttt{cbs} on top-1.}
\label{tab:twostream-ablation}
\end{table}

\subsection{Ensemble composition}

The closed-track ensemble combines the TSM variants of
Table~\ref{tab:tsm-ablation} with the $bc{=}32$ two-stream variants of
Table~\ref{tab:twostream-ablation}. With simplex-SGD weights and
per-model temperatures (fit on \texttt{val} / OOF) plus 5-crop TTA, the
ensemble reaches top-1 $0.40$ / top-5 $0.73$, well above the best single
model ($0.33$). The weak $bc{=}64$ two-stream checkpoints are kept only
as decorrelated members; the simplex assigns them near-zero weight.

\subsection{Hyperparameter table}

\begin{table}[h]
\centering
\small
\begin{tabular}{lcccc}
\toprule
\textbf{Family} & \textbf{Opt.} & \textbf{lr} & \textbf{Ep.} & \textbf{BS} \\
\midrule
CNN--LSTM              & AdamW & $3\!\cdot\!10^{-4}$ & 60 & 32 \\
Two-stream             & AdamW & $3\!\cdot\!10^{-4}$ & 60 & 32 \\
TSM                    & AdamW & $1\!\cdot\!10^{-3}$ & 30 & 32 \\
VideoMAE base (FT)     & AdamW & $5\!\cdot\!10^{-5}$ &  5 & 48 \\
TimeSformer base (FT)  & AdamW & $5\!\cdot\!10^{-5}$ &  5 & 24 \\
MViTv2 base (FT)       & AdamW & $1\!\cdot\!10^{-4}$ &  4 &  8 \\
V-JEPA 2 0.4B (LoRA)   & AdamW & $1\!\cdot\!10^{-4}$ &  5 &  4 \\
V-JEPA 2 1B (LoRA)     & AdamW & $1\!\cdot\!10^{-4}$ &  5 &  1 \\
\bottomrule
\end{tabular}
\caption{Per-family hyperparameters (best configuration). All other
values are inherited from \texttt{configs/train/default.yaml}.}
\label{tab:hparams}
\end{table}

{\small
\bibliographystyle{ieee_fullname}
\begin{thebibliography}{99}
\bibitem{Goyal17} R. Goyal \etal, ``The Something Something Video Database for Learning and Evaluating Visual Common Sense,'' \emph{ICCV}, 2017.
\bibitem{Zhang18mixup} H. Zhang \etal, ``mixup: Beyond Empirical Risk Minimization,'' \emph{ICLR}, 2018.
\bibitem{Yun19cutmix} S. Yun \etal, ``CutMix: Regularization Strategy to Train Strong Classifiers with Localizable Features,'' \emph{ICCV}, 2019.
\bibitem{Jozefowicz15} R. Jozefowicz \etal, ``An Empirical Exploration of Recurrent Network Architectures,'' \emph{ICML}, 2015.
\bibitem{Lin19tsm} J. Lin, C. Gan, S. Han, ``TSM: Temporal Shift Module for Efficient Video Understanding,'' \emph{ICCV}, 2019.
\bibitem{Tong22videomae} Z. Tong \etal, ``VideoMAE: Masked Autoencoders are Data-Efficient Learners for Self-Supervised Video Pre-Training,'' \emph{NeurIPS}, 2022.
\bibitem{Bertasius21timesformer} G. Bertasius, H. Wang, L. Torresani, ``Is Space-Time Attention All You Need for Video Understanding?,'' \emph{ICML}, 2021.
\bibitem{Li22mvitv2} Y. Li \etal, ``MViTv2: Improved Multiscale Vision Transformers for Classification and Detection,'' \emph{CVPR}, 2022.
\bibitem{Hu22lora} E. Hu \etal, ``LoRA: Low-Rank Adaptation of Large Language Models,'' \emph{ICLR}, 2022.
\bibitem{Menon20logitadj} A. Menon \etal, ``Long-Tail Learning via Logit Adjustment,'' \emph{ICLR}, 2021.
\bibitem{Guo17calibration} C. Guo \etal, ``On Calibration of Modern Neural Networks,'' \emph{ICML}, 2017.
\end{thebibliography}
}

\end{document}
